\subsubsection{Watchpoint za pronalaženje korupcije memorije}\label{gdb-watchpoints}

Tačka prekida zaustavlja program kada izvršavanje dođe do određene linije ili funkcije. \emph{Watchpoint} zaustavlja program kada se promeni određena memorijska lokacija. To je korisno kada znamo koja promenljiva dobija pogrešnu vrednost, ali ne znamo ko ju je prepisao.

Primer ima ručno napisano kopiranje stringa bez provere granice:

\codefile{watchpoint_demo.c}{c}

Prevedimo program:

\begin{minted}{bash}
$ make
$ ./watchpoint_demo petar
$ ./watchpoint_demo AAAAAAAAAAAA
\end{minted}

Drugo pokretanje može promeniti polje \texttt{is\_admin}, jer se ulaz upisuje preko kraja niza \texttt{name}. U \texttt{gdb}-u prvo zaustavljamo program posle inicijalizacije strukture, a zatim postavljamo watchpoint:

\begin{minted}{text}
(gdb) break watchpoint_demo.c:23
(gdb) run AAAAAAAAAAAA
(gdb) watch account.is_admin
(gdb) continue
\end{minted}

Kada se polje promeni, \texttt{gdb} ispisuje staru i novu vrednost i zaustavlja program na instrukciji koja je izvršila upis. Zatim proveravamo kontekst:

\begin{minted}{text}
(gdb) bt
(gdb) frame 0
(gdb) info locals
(gdb) p i
(gdb) p account
\end{minted}

Očekivani zaključak je da problem nije u kodu koji čita \texttt{is\_admin}, nego u funkciji \texttt{copy\_name}. Kada brojač \texttt{i} pređe veličinu niza, sledeći upisi prelaze u naredna polja strukture.

Hardverski watchpoint-i su ograničen resurs. Procesor obično može da prati mali broj adresa istovremeno. Ipak, kada je simptom korupcija jedne promenljive, ovo je često brže od ručnog pregleda celog toka programa.
